\chapterOrPart{Useful commands}\label{ch:usefulcommands}
This chapter contains an alphabetic list of some useful commands. All the commands that are mentioned here are included in the style file \verb|myCommands.sty|.

\section{Basic commands in \LaTeX}
There are several good sources to learn about commands in \LaTeX. See for instance the list below
\begin{itemize}
    \item \href{https://www.ntg.nl/doc/biemesderfer/ltxcrib.pdf}{\LaTeX\ command summary}
    \item \href{https://wch.github.io/latexsheet/}{\LaTeX\ cheat sheet}
    \item \href{https://en.wikibooks.org/wiki/LaTeX/Command_Glossary}{Wikibook: \LaTeX\ command glossary} 
\end{itemize}

Below we will mention some commands that might be of interest for you.

\section{definecolor}
You can define your own colours. For instance:

\verb|\definecolor{VUBorange}{cmyk}{0.00, .78, 1.00, 0.00}|

will define the colour VUBorange. This colour can now be used in your document, as is done on the title page. For instance, the command 

\verb|\textcolor{VUBorange}{This text is in VUB orange}|

will produce: \textcolor{VUBorange}{This text is in VUB orange}.

\section{ensuremath}
The command \verb|\ensuremath| switches to math-mode, if necessary. It is typically used inside \verb|\newcommand|. For instance, the command \verb|$\psat$| will produce $\psat$ but the command without the dollar-signs \verb|\psat| will produce exactly the same result. This is because the command \verb|\ensuremath| is used inside the definition of \verb|\psat| (you can verify this in the style file \verb|myCommnands.sty|) and that will force a switch to math-mode, if not yet in math-mode. This is very convenient because you are no longer obliged to explicitly activate math-mode with the \$-signs, every time you use such a command when writing.

\section{includeonly}
This command is used to reduce the compilation time of your Overleaf project. There are actually two ways to reduce compilation time:
\begin{enumerate}
    \item Reduce the amount of content in the document part. This can be achieved by disabling \verb|\include{path/filename}| commands. Just put a \% sign at the start of the line. For instance: \%\verb|\chapterOrPart{Useful packages}\label{ch:usefulpackages}
This chapter contains an alphabetic list of some useful packages. All the packages that are mentioned here are included in the style file \verb|styles/myPackages.sty|.

\section{biblatex}\label{sec:biblatex}
This package is used for citing and formatting the bibliography. It is a complete implementation of the bibliographic facilities provided by \LaTeX. BibLaTeX provides several standard citations styles, as for example: \verb|authoryear|, \verb|numeric| and \verb|alphabetic| amongst many other non-standard citation styles that are popular in different journals.

You can inspect the current options that are loaded for the BibLaTeX package in the style file \verb|styles/myPackages.sty|. Some common styles and sorting options are shown in \cref{tab:biblatex}. 

\begin{table}\centering
    \caption{Some examples of style and sorting options in BibLaTeX}\label{tab:biblatex}
        \begin{tabular}{lllll}
        \toprule
        \verb|style|& \verb|sorting| & \verb|\cite| example  & cite number         & order in bibliography  \\[+1.5ex]
        authoryear  & nyt            & Rees, 2017 &                     & by name, year, title \\
        authoryear  & none           & Rees, 2017 &                     & by cite order \\
        numeric     & nyt            & [8]        &  as in bibliography & by name, year, title \\
        numeric     & none           & [1]        &  as order in text   & by cite order  \\
        alphabetic  & nyt            & [Ree17]    &                     & by name, year, title \\
        alphabetic  & none           & [Ree17]    &                     & by cite order  \\
        \bottomrule
        \end{tabular}
\end{table}

\begin{table}\centering
    \caption{Some examples of other cite commands}\label{tab:cites}
        \begin{tabular}{rlllll}
        \toprule
            \verb|style|: 	            &	current	style           &	authoryear	&	numeric	    &	alphabetic	\\[+1.5ex]
            \verb|\cite{Rees2017}| 	    &	\cite{Rees2017}	        &	Rees, 2017	&	[1]	        &	[Ree17]	    \\
            \verb|\textcite{Rees2017}| 	&	\textcite{Rees2017}	    &	Rees (2017)	&	Rees [1]	&	Rees [Ree17]\\
            \verb|\parencite{Rees2017}| &	\parencite{Rees2017}	&	(Rees, 2017)&	[1]	        &	[Ree17]	    \\
            \verb|\citeauthor{Rees2017}|&	\citeauthor{Rees2017}	&	Rees	    &	Rees	    &	Rees	    \\
        \bottomrule
        \end{tabular}
\end{table}

Is is good practice to keep your bibliographic references in a separate file, a so-called bibliographic information file, recognised by the extension \verb|.bib|. In this template, this file can be found here: \verb|bib/myBibliography.bib|.

This file is imported with the command \verb|\addbibresource{bib/myBibliography.bib}| in the style file \verb|styles/myPackages.sty|. 

It is convenient and recommended to use separate software to collect and manage your references. A well-known free tool to manage your references is \href{https://www.mendeley.com/}{Mendeley Reference Manager}. It works well together with \href{https://www.mendeley.com/reference-management/web-importer}{Mendeley Web Importer}. This is a browser extension to import reference material that you are watching in your browser into your Mendeley database.  Other well-known free tools are \href{https://www.zotero.org/}{Zotero} and \href{https://www.jabref.org/}{JabRef}, but there exist many others.

Every reference in the bib file must have a unique identifier, a so-called a \verb|BibTeX key|. You can choose this key yourself. An often-used approach to construct this key is to combine the name of the first author and the year of publishing. You can then refer to a BibTeX key with commands such as \verb|\cite|, \verb|\parencite|, \verb|\textcite| or \verb|\citeauthor|. 

See \cref{tab:cites} for some examples of these cite commands. The second column in this table shows the results for the style that is currently loaded in BibLaTeX. \verb|\cite| is used in sentences where the reference is part of the sentence, and the year is there to identify what is being cited. \verb|\parencite| is used to put the citation between parentheses. It is used when the sentence is already complete, and the citation is being added in support. \verb|\textcite| is used in sentences where the name of the author is part of the sentence, and the year is there to identify what is being cited. With \verb|\citeauthor| only the name of the author is shown without a hyperlink to the reference in the bibliography. 

\textcite{Rees2017} made a good overview of the available commands and options in  \href{https://tug.ctan.org/info/biblatex-cheatsheet/biblatex-cheatsheet.pdf}{this cheat sheet}.

More info at \href{https://ctan.org/pkg/biblatex?lang=en}{CTAN: Package BibLaTeX} and at \href{https://www.overleaf.com/learn/latex/Articles/Getting_started_with_BibLaTeX}{Overleaf: getting started with BibLaTeX}.

\section{cancel}\label{sec:cancel}
With this package you can easily show which terms are canceled in equations. For instance, the command \verb|$\cancel{x}$| will produce $\cancel{x}$ and \verb|$\cancelto{0}{z}$| will produce $\cancelto{0}{z}$.

Both commands are also demonstrated in an equation below.
\begin{frame}{With the package \texttt{cancel}, you can cancel terms in equations}
    \begin{align}
    % the command \intertext is handy to put text between equations in the align environment.
    \intertext{For example, for a stationary \dutch{stilstaand} process:}
    &\Delta E_\CV=\Delta U_\CV + \cancel{\Delta\up{KE}_\CV}\quad +\quad\cancelto{\SI{0}{J}}{\Delta\up{PE}_\CV} \\
    \intertext{and the first law for a closed and stationary system simplifies to}
    &\Delta E_\CV=\Delta U_\CV =Q+W\units{J}.
    \end{align}
\end{frame}

More info at \href{https://ctan.org/pkg/cancel?lang=en}{CTAN: Package cancel}.

\section{chemformula}\label{sec:chemformula}
This package allows you to write chemical formulas in a simple way. Consider for instance the formula $\ch{H2O}$ for water. In the table below you will find some bad and good ways to write it in \LaTeX.

\begin{frame}[fragile]{Some good and bad ways to write a chemical formula}\smallbeamer
% option [fragile] needed when using verbatim on slide with \verb|something|
    \begin{center}
        \begin{tabular}{lll}
        \toprule
        \LaTeX\ command & result & remark \\[+1.5ex]
        \verb|H2O|  & H2O &  not correct, 2 not in subscript \\
        \verb|H_2O| & error & underscore not allowed outside math-mode \\
        \verb|$H_2O$|     & $H_2O$& bad practice to write chemicals in italics\\
        \verb|$\mathrm{H_2O}$|   & $\mathrm{H_2O}$ & correct, but heavy syntax\\
        \verb|\ch{H2O}|   & \ch{H2O} & correct and easy syntax\\
        \bottomrule
        \end{tabular}
    \end{center}
\end{frame}

Writing chemical reactions is also possible. For instance:

\verb|\ch{H2O + CO3^2- <=> OH- + HCO3-}|

produces: \ch{H2O + CO3^2- <=> OH- + HCO3-}

More info at \href{https://ctan.org/pkg/chemformula?lang=EN}{CTAN: Package chemformula}.

\section{cleveref}\label{sec:cleveref}
The \texttt{cleveref} package enhances \LaTeX's cross-referencing features. It automatically formats cross-references depending on what they refer to (chapter, section, table, figure, equation,\dots). The normal procedure in \LaTeX\ is to use the \verb|\ref| command preceded with the type of object that is being referred to. With the \verb|\cref| command, this is no longer necessary because the type of object will be automatically detected. This will make your referring more consistent. In \cref{tab:cref} you will see the difference. For instance in the case of figures, you don't have to explicitly write \verb|fig.| in front of every reference to a figure. For references at the beginning of a sentence, the \verb|\Cref| command must be used. 

\begin{table}\centering
    \caption{Some examples of the use of \texttt{\textbackslash cref} and \texttt{\textbackslash Cref}}\label{tab:cref}
        \begin{tabular}{lll}
        \toprule
        command                         & old syntax with \verb|\ref|           & result \\[+1ex]
        \verb|\cref{ch:conclusions}|    & \verb|chapter~\ref{ch:conclusions}|   & \cref{ch:conclusions} \\
        \verb|\cref{sec:siunitx}|       & \verb|section~\ref{sec:siunitx}|      & \cref{sec:siunitx} \\
        \verb|\cref{fig:psat}|          & \verb|fig.~\ref{fig:psat}|            & \cref{fig:psat} \\
        \verb|\cref{tab:mathsymbols}|   & \verb|table~\ref{tab:mathsymbols}|    & \cref{tab:mathsymbols} \\
        \verb|\cref{eq:newton}|         & \verb|eq.~(\ref{eq:newton})|          & \cref{eq:newton} \\
        \verb|\cref{list:factorial}|    & \verb|listing~(\ref{list:factorial})| & \cref{list:factorial} \\
        \midrule
        %\verb|\Cref| commands          &                                       &  \\[+1.5ex]
        \verb|\Cref{ch:conclusions}|    & \verb|Chapter~\ref{ch:conclusions}|   & \Cref{ch:conclusions} \\
        \verb|\Cref{sec:siunitx}|       & \verb|Section~\ref{sec:siunitx}|      & \Cref{sec:siunitx} \\
        \verb|\Cref{fig:psat}|          & \verb|Figure~\ref{fig:psat}|          & \Cref{fig:psat} \\
        \verb|\Cref{tab:mathsymbols}|   & \verb|Table~\ref{tab:mathsymbols}|    & \Cref{tab:mathsymbols} \\
        \verb|\Cref{eq:newton}|         & \verb|Equation~(\ref{eq:newton})|     & \Cref{eq:newton} \\
        \verb|\Cref{list:factorial}|    & \verb|Listing~(\ref{list:factorial})| & \Cref{list:factorial} \\
        \bottomrule
        \end{tabular}
\end{table}

The format for each cross-reference type (chapter, section, figure, table and equation) can be customised in the style file \verb|myPackages.sty|. 

More info at \href{https://ctan.org/pkg/cleveref}{CTAN: Package cleveref}.

\section{hyperref}\label{sec:hyperref}
With this package you can add hyperlinks in your document. For instance, the command:

\verb|\href{https://canvas.vub.be/}{Canvas learning platform}|

will produce this clickable link: \href{https://canvas.vub.be/}{Canvas learning platform}.

More info can be found in \href{https://ctan.org/pkg/hyperref}{CTAN: Package hyperref}.

\section{Listings}\label{sec:listing}
Sometimes, you want to add source code to your document. Proper typesetting will make your code easier to read. General info about typesetting source code in \LaTeX\ can be found in the \href{https://www.overleaf.com/learn/latex/Code_listing}{Overleaf help pages}. You will read that the package \texttt{Listings} is well suited for this propose. The package provides the \texttt{lstlisting} environment as shown in the "Hello world!" example in \cref{list:hello}. As can be seen in the source code of this example, you can specify the language your code is written in, add a caption and a label for cross-referencing.

Many programming languages can be typeset with this package. The names of the recognised languages can be consulted in \href{https://mirror.lyrahosting.com/CTAN/macros/latex/contrib/listings/listings.pdf}{Tabel 1 in the package documentation}. 

\begin{verbatim}
    \begin{lstlisting}[language=Matlab,label={list:hello},
    caption={Hello world in Matlab}]
        disp('Hello, World!');    
    \end{lstlisting}
\end{verbatim}

The result is shown in \cref{list:hello}.

\begin{lstlisting}[language=Matlab,label={list:hello},caption={Hello world in Matlab}]
    disp('Hello world!');  
\end{lstlisting}

Another example of typesetting of source code, this time in Python, is shown in \cref{list:factorial}. Many typesetting properties (such as: font, font size, indents, colours, \dots) can be specified. The package settings in this template can be consulted and adapted in the style file \texttt{styles/myListings.sty}.

\begin{frame}[fragile]{Listings of code can be easily formatted and displayed}
    \begin{lstlisting}[language=Python, caption={Python code to calculate factorial, adapted from \href{https://www.programiz.com/python-programming/examples/factorial-recursion}{programiz}},label={list:factorial}]
    # Factorial of a number via recursion
    def recur_factorial(n):
       if n == 1:
           return n
       else:
           return n*recur_factorial(n-1)
    
    num = int(input("Enter a number: "))
    if num < 0: # check if the number is negative
       print("Sorry, factorial does not exist for negative numbers")
    elif num == 0: # check if the number is zero
       print("The factorial of 0 is 1")
    else:
       print("The factorial of", num, "is", recur_factorial(num))
    \end{lstlisting}
\end{frame}

It is also possible to import code from a file (that you maybe are still working on) with the command \verb|\lstinputlisting{sourcecode_filename.py}|.

More info at \href{https://ctan.org/pkg/listings?lang=en}{CTAN: Package listings}.

\section{longtable}
With this package you can make tables that are larger than one page and continue to the next page(s).

More info can be found in \href{https://ctan.org/pkg/longtable?lang=en}{CTAN: Package longtable}.

\section{pgfplots}
With this package you can make high-quality plots in normal or logarithmic scaling. They are typically drawn in the \verb|axes| environment and plots are added with the command \verb|\addplot|. See for instance \cref{fig:analyticalfunctions} with some plots of analytical functions and \cref{fig:thermalpower} with a plot of external data from a file. See \href{http://pgfplots.sourceforge.net/gallery.html}{this gallery for many more examples}.

More info can be found in \href{https://ctan.org/pkg/pgfplots?lang=en}{CTAN: Package pgfplots}.

\section{siunitx}\label{sec:siunitx}
This package will neatly and correctly format your numbers and units. You will find some examples of commands and their results in \cref{tab:siunitx}.

\begin{frame}[fragile]{With the package \texttt{siuntix}, you can format numbers and units}\scriptsizebeamer
    \begin{table}\centering
        \caption{Some examples of commands available in \texttt{siunitx}}
           % \begin{tabular}{p{8cm}l}
            \begin{tabular}{ll}
                \toprule
                \LaTeX\ command & result \\[+1ex]
                \verb|\num{3.1E3}| & \num{3.1E3} \\ % to correctly format numbers
                \verb|\numlist{50;70;90}| & \numlist{50;70;90}\\ % list of numbers
                \verb|\SI{50}{\degree}| & \SI{50}{\degree}\\ % to format an angle
                \verb|\SI{31}{\celsius}| & \SI{31}{\celsius}\\ % to format an angle
                \verb|\ang{50;5;33}| & \ang{50;5;33}\\ % to format an angle
                \verb|\SI{9.81}{m/s^2}| & \SI{9.81}{m/s^2} \\ % unit symbols will always produce inclined fraction line in units
                \verb|\SI{9.81}{\meter\per\second\squared}| & \SI{9.81}{\meter\per\second\squared} \\ % unit names will format fraction as specified in style file
                \verb|\SI{9.81}{\meter\per\square\second}| & \SI{9.81}{\meter\per\square\second} \\ % unit names will format fraction as specified in style file
                \verb|\SI{9.81}[per-mode=symbol]{\meter\per\square\second}| & \SI[per-mode=symbol]{9.81}{\meter\per\square\second} \\ % inclined fraction line in units
                \verb|\si{\joule\per\second}| & \si{\joule\per\second} \\ % to format units only
                \verb|\SI{5.0}{\joule\per\second}| & \SI{5.0}{\joule\per\second} \\
                \verb|\si[per-mode=fraction]{\joule\per\second}| & \si{\joule\per\second} \\ % horizontal fraction line in units
                \verb|\si[per-mode=symbol]{\joule\per\second}| & \si[per-mode=symbol]{\joule\per\second} \\ % inclined fraction line in units
                \verb|\si{g_{vapour}\per kg_{dry air}}| & \si{g_{vapour}\per kg_{dry air}}\\
                \verb|\SI{4}{\euro\per\kg}| & \SI{4}{\euro\per\kg} \\
                \verb|\SI{2E-4}{\mol\per\cm\squared\per\second}| & \SI{2E-4}{\mol\per\cm\squared\per\second} \\[+1.5ex]
                \bottomrule
            \end{tabular}\label{tab:siunitx}
    \end{table}
\end{frame}

The package also introduces a column type \texttt{S} to align numbers in tables at the decimal separator. Inspect the source code of \cref{tab:tabledemo} to see the application of these column types. Column \texttt{unit} is centred and uses the command \verb|\si{}| to write the units. The columns indicating the initial end final state are of type \texttt{S}.

\begin{frame}{Formatting numbers and units in tables with \texttt{siunitx}}
    \begin{table}[b]\centering
    	\caption{Example of a property table with type \texttt{S} columns to align numbers at decimal separator}
    			\only<presentation>{\footnotesize}
    	\begin{tabular}{ccScSl} % column types: centre, centre, S, centre, S, left align
    		% column type s for units, column type S for values aligned at decimal separator
                \toprule
                 property & unit            & {initial state}       & process & {end state}         & comment   \\[+1.5ex]
                 $m$      & \si{\kilogram}  &  1.0                  &  $=$    &      1.0            & closed      \\
                 $T$      & \si{\celsius}   &  25.0                 &  $=$    &     25.0            & isotherm    \\
                 $\vol$   & \si{m^3}        &  1.00                 &  $>$    &     0.83            & compression \\
                \bottomrule
    	\end{tabular}\label{tab:tabledemo}
    \end{table}
\end{frame}

More info can be found in \href{https://ctan.org/pkg/siunitx?lang=en}{CTAN: Package siunitx}.


\section{TikZ}\label{sec:tikz}
\verb|TikZ| is a powerful drawing tool inside \LaTeX. Here are useful links to learn \verb|TikZ|
\begin{itemize}
    \item \href{https://www.overleaf.com/learn/latex/TikZ_package}{Overleaf documentation on TikZ}
    \item \href{https://cremeronline.com/LaTeX/minimaltikz.pdf}{A very minimal introduction to TikZ}
    \item \href{https://www.math.uni-leipzig.de/~hellmund/LaTeX/pgf-tut.pdf}{PGF/TikZ - Graphics for \LaTeX} 
    \item \href{http://siam.cs.kuleuven.be/sites/siam.cs.kuleuven.be/files/events/tikz_workshop/tikz_introduction.pdf}{A very short introduction to TikZ} 
    
\end{itemize}

You can for instance make simple inline drawings with the command \verb|tikz{}|. For instance the command

\verb|\tikz{\draw[blue] (0,0) rectangle (3cm,2mm)}|

will draw this blue rectangle \tikz{\fill[blue] (0,0) rectangle (3cm,2mm)}.

More complex drawings are typically made within the \verb|tikzpicture| environment. See \cref{ch:drawing} in this document for some examples. Many more examples can be found on \href{http://www.texample.net/tikz/examples/}{TeXample.net}.

More info can be found in \href{https://ctan.org/pkg/tikz?lang=en}{CTAN: Package TikZ}.

\section{todonotes}\label{sec:todonotes}
With the command \verb|\todo| you can insert notes in your document about things that needs to be
done later. For instance the command \verb|\todo{example of a margin note}| renders to the note in the margin you see in the pdf of this document.
\todo{example of a margin note}

The command \verb|\missingfigure{Search for better picture.}|

will insert a placeholder for a picture you still need to add as shown below.
\missingfigure{Search for better picture.}

The command \verb|\listoftodos| will add a list of all todo's. See last page in this document for such a list.

More info can be found in \href{https://ctan.org/pkg/todonotes?lang=en}{CTAN: Package todonotes}.

\mode<all> 
% Required because of the \mode* command at the beginning of this file.
% See beameruserfile.pdf section 21.3 for more info|, will no longer compile the \verb|usefulPackages| file. As a result:
        \begin{itemize}
            \item The content of \verb|usefulPackages| will no longer be known to \LaTeX, and thus will no longer show up in the output pdf
            \item The content of \verb|usefulPackages| will no longer show up in the table of contents
            \item Labels in \verb|usefulPackages| will no longer be known to \LaTeX, so cross-references to such labels, from anywhere in your document, will result in broken links, indicated by a double question mark: ??
        \end{itemize}
    \item Use the command \verb|\includeonly{path/filename}| in the preamble. But beware, you have to do proceed as follows:
    \begin{enumerate}
        \item First comment out \%\verb|\includeonly{path/filename}| in the preamble and compile your entire project. This will allow \LaTeX\ to detect the document structure (chapters, sections, labels, \dots) in all the included files.
        \item After successful compilation, you add the file you want to work in today. Let's assume it is the \verb|usefulPackages| file. In that case, you add \\ \verb|\includeonly{appendices/usefulPackages.tex}| to your preamble, not to the document part!
        \item Then you-recompile your project. As a result:
    \end{enumerate}
      \begin{itemize}
            \item Only the content of \verb|usefulPackages| will be recompiled by \LaTeX. The other included files are not recompiled. This will speed up compilation! The content of all the other compiled files is however still known to \LaTeX\ from the previous full compile. 
            \item Only the content of \verb|usefulPackages| will show up in the output pdf.
            \item The content of all included files will show up in the table of contents, because \LaTeX\ remembers the document structure from the first entire compile.
            \item All the labels in all compiled files are known to \LaTeX, so cross-references to such labels, from anywhere in your document, will work.
        \end{itemize}
\end{enumerate}

\section{intertext}
The command \verb|\intertext| is useful to write lines of text between equations. It is applied in the \verb|align| environment. You can then use one single \verb|align| environment, instead of a series of separate \verb|equation| environments. Inspect the source code in the example below. This approach is also very well suited for writing equations on slides.
\begin{frame}{With the \texttt{intertext}-command, text can easily be placed between aligned equations}
    \begin{align}
        \intertext{For an adiabatic CV, the entropy balance is} 
        &\dod{S_\CV}{t}=\sum_\up{in}\mdot_i\cdot s_i-\sum_\up{out}\mdot_j\cdot s_j+\dot{S}_\up{gen}\\
        \intertext{for a closed CV with $k$ heat flows, it takes the following form}
        &\dod{S_\CV}{t}=\sum\frac{\qdot_k}{T_k}+\dot{S}_\up{gen}\\
        \intertext{and for an isolated CV it reduces to}
        &\dod{S_\CV}{t}=\dot{S}_\up{gen}\geq 0\qquad\text{and thus $S$ reaches a maximum at equilibrium.}
    \end{align} 
\end{frame}

\section{newcommand}
With the command \verb|\newcommand| you can define your own commands. That can be useful to simplify repetitive and/or complex formatting. See the \href{https://www.overleaf.com/learn/latex/Commands#Defining_a_new_command}{Overleaf documentation on newcommand} for more info and the examples below.

\begin{frame}[fragile]{Some examples of self-defined commands}\small % size for report
    \begin{center}\scriptsizebeamer % reduce size to fit on slide
        \begin{tabular}{l}
        \toprule
        \verb|\newcommand{\vel}[0]{\ensuremath{V}}|\\
        example: \verb|\vel= \SI{2}{m/s}|\\
        result: \vel= \SI{2}{m/s} \\[+1ex]
        
        \verb|\newcommand{\psat}[0]{\ensuremath{\pres_\up{sat}}}|\\
        example: \verb|\psat = \SI{3.2}{\kilo\pascal}|\\
        result: \psat = \SI{3.2}{\kilo\pascal}\\[+1ex]
    
        % shorter version of \mathrm
        % \verb|\newcommand\up[1]{\mathrm{#1}}|\\ 
        % example: \verb|$T_\up{air} = \SI{25}{\celsius}$|\\
        % result: $T_\up{air} = \SI{25}{\celsius}$\\[+1ex]
       
        % to put units behind an equation 
        \verb|\newcommand{\units}[1]{\ensuremath{\quad\up{\left[#1\right]}}}|\\ 
        example: \verb|$T \ud s = \delta q + \delta d \units{J/kg}$|\\
        result: $T \ud s = \delta q + \delta d \units{J/kg}$\\[+1ex]
        
        % to put units behind an equation
        \verb|\newcommand{\trans}[1]{[{\footnotesize Dutch:} {#1}]}|\\
        example: \verb|The flow is steady \trans{stationair}.|\\
        result: The flow is steady \dutch{stationair}.\\[+1ex] 
        
        % specific heat capacity of water
        \verb|\newcommand{\cvwatervalue}[0]{\SI{4.18}{\kilo\joule\per\kilogram\per\kelvin}}|\\
        example: \verb|\cvwatervalue|\\
        result: \cvwatervalue\\[+1ex] 
        \bottomrule
        \end{tabular}
    \end{center}
\end{frame}

\section{newtoggle}

You can define your own boolean variables called \verb|toggles|. These toggles can be used to make certain decisions.

For instance, the command \verb|\newtoggle{showResults}| defines the toggle \verb|showResults|.

You can set this toggle to TRUE with the command \verb|\toggletrue{showResults}| and to FALSE with \verb|\togglefalse{showResults}|.

We can now use this toggle to hide or show some text, e.g.

the command

\verb|$1+1=3$ \iftoggle{showResults}{for very large values of 1.}{when ?}|

will produce:

\togglefalse{showResults}
$1+1 = 3$ \iftoggle{showResults}{for very large values of 1.}{when ?} (when \verb|showResults| is set to FALSE)

\toggletrue{showResults}
$1+1 = 3$ \iftoggle{showResults}{for very large values of 1.}{when ?} (when \verb|showResults| is set to TRUE)

\section{textcolor}
With this command you can colour text. For instance the command

\verb|\textcolor{blue}{This text is coloured}|

will produce: \textcolor{blue}{This text is coloured}.


\mode<all> 
% Required because of the \mode* command at the beginning of this file.

% See beameruserfile.pdf section 21.3 for more info